% SOURCE_AUTHORITY: sga5_fr_workpass.tex, lines 573--578 (LF-normalized unit).
\subsection*{Corolário 1.15}
Seja $f : X \to Y$ um morfismo quase-finito e separado de esquemas noetherianos. Seja $K_Y \in \ob \D^{+}(Y)$, e ponhamos, como em 1.12, $K_X = \R^{!}f(K_Y)$, $\underline{D}_X = \underline{D}_{K_X}$, $\underline{D}_Y = \underline{D}_{K_Y}$. Então, se $K_Y$ é dualizante, $K_X$ é dualizante.

\emph{Demonstração.} Já sabemos (1.5) que, se $K_Y$ tem dimensão quase-injetiva finita, o mesmo ocorre com $K_X$. Demonstremos que, se $K_Y$ satisfaz a condição (ii) (resp.~(iii)) de (1.7), $K_X$ também a satisfaz. Segundo o \textquotedblleft teorema principal\textquotedblright{} (EGA~IV~8.12.6), $f$ se fatora como $f'i$, onde $i$ é uma imersão aberta e $f'$ um morfismo finito. Como $\R^{!}f \simeq \R^{!}i \, \R^{!}f'$, deve-se considerar separadamente o caso de uma imersão aberta e o de um morfismo finito. Se $f$ é uma imersão aberta, a afirmação é trivial, pois um feixe construtível sobre $X$ é induzido por um feixe construtível sobre $Y$, e a formação dos $\R\,\underline{\Hom}$ comuta com a restrição a um aberto. Suponhamos, pois, que $f$ é finito. Sabemos então, por (1.13), que, se $K_Y$ satisfaz (iii), $K_X$ também o satisfaz (levando em conta o teorema de finitude SGAA~IX~2.14). Resta demonstrar que, se $K_Y$ satisfaz (ii), $K_X$ também o satisfaz. Seja $F \in \ob \D_c^{-}(X)$; demonstremos que $\underline{D}_X(F) \in \ob \D_c^{+}(X)$. Pelo teorema de finitude, tem-se $f_{*}(F) \in \ob \D_c^{-}(Y)$ e, portanto, por hipótese, $\underline{D}_Y f_{*}(F) \in \ob \D_c^{+}(Y)$. Mas $\underline{D}_Y f_{*}(F) \simeq f_{*}\underline{D}_X(F)$ (1.12~a)(i)), e a conclusão segue de (1.14~b)).
